\documentclass[11pt,a4paper]{article}
\usepackage[margin=2.2cm]{geometry}
\usepackage{booktabs}
\usepackage{graphicx}
\usepackage{amsmath}
\usepackage{xcolor}
\usepackage{enumitem}
\usepackage[font=small]{caption}

\setlength{\parskip}{0.4em}
\setlength{\parindent}{0em}

\usepackage{graphicx} % Required for inserting images
\usepackage{enumitem}
\usepackage[colorlinks=true, linkcolor=black, urlcolor=black, citecolor=black]{hyperref}
\usepackage{float}
\usepackage{array,booktabs}
\usepackage[most]{tcolorbox}
\usepackage{amssymb}
\usepackage[numbers]{natbib}

% =============== BOX GRIGIO ===================
\tcbset{
  tabbox/.style={
    enhanced,
    colback=black!7,   % grigio chiaro
    colframe=black!7,  % stesso colore per nascondere il bordo
    boxrule=0pt,
    arc=3mm,           % angoli arrotondati
    left=8mm,right=8mm,top=6mm,bottom=6mm
  }
}

\definecolor{rosepurple}{RGB}{255,90,255}
\definecolor{orange}{RGB}{255,50,50}
\definecolor{blue}{RGB}{90, 90, 255}
\definecolor{codegreen}{rgb}{0,0.6,0}
\definecolor{codegray}{rgb}{0.5,0.5,0.5}
\definecolor{codepurple}{rgb}{0.58,0,0.82}
\definecolor{backcolour}{rgb}{0.95,0.95,0.92}
\DeclareMathOperator*{\argmax}{arg\,max}
\DeclareMathOperator*{\argmin}{arg\,min}

\lstdefinestyle{mystyle}{
    backgroundcolor=\color{backcolour},
    commentstyle=\color{codegreen},
    keywordstyle=\color{magenta},
    numberstyle=\tiny\color{codegray},
    stringstyle=\color{codepurple},
    basicstyle=\footnotesize\ttfamily,
    breakatwhitespace=false,
    breaklines=true,
    captionpos=b,
    keepspaces=true,
    numbers=left,
    numbersep=5pt,
    showspaces=false,
    showstringspaces=false,
    showtabs=false,
    tabsize=2
}

\lstset{style=mystyle}
\begin{document}
\begin{titlepage}
    \centering
    \vspace*{4cm}

    \includegraphics[width=1\textwidth]{jetbrains-logo.png}\par
    \vspace{2cm}

    {\Huge \bfseries Stem Agent: Self-Specializing AI Agent Through Guided Evolutionary Configuration Search\par}
    \vspace{1cm}
    {\Large Andrea Gennari\par}

    \vfill
\end{titlepage}

\tableofcontents


% =============================================================
% DAY 1: Write Section 1 + Section 2 draft
% =============================================================
\newpage

\section{Introduction and Approach}

A stem cell begins life in a generic, pluripotent state and differentiates into a specialist by reading environmental signals.
We apply the same principle to AI agents.
A \emph{Stem Agent} starts with a minimal, generic configuration for a target problem class and iteratively specializes itself: it mutates its own configuration, measures the effect on a held-out performance signal, and accepts or rejects the change subject to a safety invariant that prevents catastrophic regression.
The result is an agent that earned its specialization through evidence, not intuition.

We choose \textbf{Python bug repair} as the target domain for three concrete reasons.
First, evaluation is perfectly objective: a test suite either passes or it does not, with no scoring rubric or human judgment required.
Second, the domain is directly relevant to JetBrains, whose core business is developer productivity tooling.
Third, bug repair has structured failure modes such as off-by-one boundary shifts, inverted conditions, wrong variable references, and missing guard clauses, which reward specialization in a measurable and reproducible way.

The architecture consists of four sequential phases.
\textbf{Phase~0 (Baseline)} establishes a performance floor by running a vanilla agent on all 90 tasks and using its per-task pass/fail results to assign four frozen evaluation splits.
\textbf{Phase~1 (Sensing)} analyzes failures on a calibration split using an LLM to extract structured failure patterns and to suggest specialization directions; this produces a static \texttt{SensorReport}.
\textbf{Phase~2 (Differentiation)} runs an evolutionary mutation-selection loop: each iteration generates a \emph{challenger} configuration by applying one or two mutation operators to the current \emph{champion}, evaluates it on the dev set, checks it against a vital signs safeguard, and promotes it if it clears both thresholds.
\textbf{Phase~3 (Final Evaluation)} scores the final champion, together with all ablation variants, on the held-out test set that was never touched during search.

The inner agent's ``genome'' is an \texttt{AgentConfig} dataclass that covers four mutation axes: (a)~\emph{prompt} (system message, few-shot examples, chain-of-thought instruction, output format); (b)~\emph{tools} (subset of 7 available tools, tool-use strategy); (c)~\emph{strategy} (\texttt{direct}, \texttt{cot}, \texttt{react}, \texttt{plan\_execute}); and (d)~\emph{inference parameters} (temperature, model).
Twelve explicit mutation operators act on these axes (Section~\ref{sec:differentiation}).

The key design decision is the \textbf{vital signs invariant}.
After Phase~0, a small set of tasks that the vanilla baseline solves are designated \emph{vital signs}.
Any challenger configuration that fails on any vital-sign task is rejected regardless of its dev-set score.
This prevents catastrophic specialization: the agent becoming very good at hard cases while regressing on easy ones.
No existing prompt-optimization or agent-search framework (DSPy, OPRO, Promptbreeder, ADAS) provides an analogous formal regression-safety mechanism.

\textbf{Related work.}
The closest prior work falls into two categories: prompt-optimization frameworks and agent-design search systems.

\textbf{Prompt-optimization} frameworks treat the agent's prompt as the sole axis of change.
DSPy~\cite{khattab2023dspy} compiles few-shot prompts via bootstrapped examples but fixes the tool set, reasoning strategy, and model; it provides no mechanism to detect regression on previously-solved tasks.
OPRO~\cite{yang2023opro} iteratively rewrites system prompts using a meta-prompt as a gradient substitute; it optimizes a single aggregate metric with no guard against catastrophic forgetting of easy instances.
Promptbreeder~\cite{fernando2023promptbreeder} applies evolutionary operators to prompt strings but restricts mutation to the text axis alone and does not model multi-axis agent configurations.

\textbf{Agent-design search} expands the search space beyond prompts.
ADAS~\cite{hu2024adas} uses an LLM to propose and evaluate novel agent designs expressed as Python code; it covers a larger open-ended space but neither tracks per-task regression across iterations nor includes a formal regression-safety invariant.

\textbf{Our contribution} differs on two axes.
First, we define a \emph{multi-axis genome} covering prompt, tools, reasoning strategy, and inference parameters, and apply typed mutation operators to each axis, whereas all prior work mutates prompt text alone or uses unconstrained code generation.
Second, we introduce the \textbf{vital signs invariant}: any challenger configuration that fails on any task the vanilla baseline solved is immediately rejected, regardless of its aggregate development-set score.
To our knowledge, no prior prompt-optimization or agent-search framework formalizes an analogous regression-safety mechanism~\cite{khattab2023dspy,yang2023opro,fernando2023promptbreeder,hu2024adas}.

\section{Benchmark and Experimental Setup}

\subsection{Benchmark}

The benchmark comprises \textbf{90 self-contained Python bug-repair tasks}.
Each task is a directory containing: \texttt{buggy.py} (the file with the injected bug), \texttt{fixed.py} (the ground-truth reference fix, never shown to agents), \texttt{test\_suite.py} (a \texttt{pytest} suite that fails on the buggy file and passes on the fixed file), and \texttt{metadata.json} (category, difficulty, source, description).

\textbf{Source.}
We initially attempted to extract tasks from BugsInPy~\cite{widyasari2020}, a dataset of real bugs from popular Python projects.
Repository-level dependencies and complex test environments made self-contained extraction infeasible within the time budget.
We therefore constructed all 90 tasks via \emph{mutation testing} on Exercism Python exercises: correct reference solutions were taken and a single-point bug was injected using known mutation operators (off-by-one boundary shift, wrong relational operator, inverted boolean condition, wrong variable reference, missing guard clause, wrong API argument, boundary index error).
All tasks carry \texttt{"source": "mutation\_exercism"} in their metadata.

\textbf{Category distribution.}
Tasks span 7 bug categories (Table~\ref{tab:benchmark}).
Difficulty is labeled \texttt{easy} (single-token or single-operator change; 66 tasks) or \texttt{medium} (multi-line or structural change; 24 tasks).

\begin{table}[h]
\centering
\caption{Benchmark category distribution (90 tasks).}
\label{tab:benchmark}
\begin{tabular}{lcc}
\toprule
Category & Count & \% \\
\midrule
\texttt{logic\_error}    & 18 & 20.0 \\
\texttt{off\_by\_one}    & 15 & 16.7 \\
\texttt{missing\_check}  & 15 & 16.7 \\
\texttt{wrong\_variable} & 12 & 13.3 \\
\texttt{api\_misuse}     & 10 & 11.1 \\
\texttt{type\_error}     & 10 & 11.1 \\
\texttt{boundary}        & 10 & 11.1 \\
\bottomrule
\end{tabular}
\end{table}

\textbf{Validation.}
Every task was validated automatically: \texttt{buggy.py} was run as the candidate solution and \emph{must fail} at least one test; \texttt{fixed.py} was run as the candidate solution and \emph{must pass} all tests.
Tasks that did not satisfy this invariant, because the bugs were too subtle for the provided tests to catch, were redesigned with targeted test cases that expose the fault.
All 90 tasks passed validation before being included.

\textbf{Test-source leakage and prompt design.}
A manual inspection of 10 randomly sampled tasks during construction revealed a
test-leakage failure mode: in the originally generated test suites, ultra-specific
assertions of the form \texttt{assert f(x) == hardcoded\_value} produced
tracebacks that effectively spelled out the correct return value, and several
tests defined a local \texttt{\_ref} function that re-implemented the intended
algorithm. Both patterns directly compress the dynamic range of any agent
comparison, since a competent LLM can solve the task by pattern-matching the
test source rather than by reasoning about the buggy code. This is the textbook
``synthetic data too easy'' failure mode.

We address this in two ways. First, 12 tasks (\texttt{task\_004}, \texttt{task\_005},
\texttt{task\_008}, \texttt{task\_011}, \texttt{task\_012}, \texttt{task\_015},
\texttt{task\_021}, \texttt{task\_027}, \texttt{task\_042}, \texttt{task\_055},
\texttt{task\_057}, \texttt{task\_058}) where the original tests most directly
revealed the fix were rewritten with hardcoded behavioral assertions and minimal
spoilers; the remaining 78 retain behavioral or property-based assertions from
the initial generation, several of which still embed reference implementations
in helper functions. The full rewrite of the 12 tasks is reproducible via
\texttt{benchmark/fix\_tasks.py}.

Second, and more importantly, the agent prompt does \emph{not} include the test
source. The agent receives only \texttt{buggy.py}; the test suite is never shown
verbatim. Tests remain the binary oracle for scoring (pass/fail) and are
indirectly accessible to ReAct-style agents through the \texttt{run\_tests} and
\texttt{get\_traceback} tools, which return the failing pytest output but not the
test source code. This design choice both removes the leakage from any
unrewritten \texttt{\_ref} helpers and matches a more realistic bug-fixing
scenario, where a developer typically has a buggy file plus a failure report
rather than the full test suite verbatim. A regression test in
\texttt{tests/test\_agent.py} (\texttt{test\_test\_suite\_code\_never\_in\_prompt})
locks in this invariant.

We re-validated all 90 tasks after these changes: \texttt{buggy.py} fails at
least one test and \texttt{fixed.py} passes every test on each task.

\textbf{Acknowledged benchmark limitations.} The benchmark is synthetic: every
task is a single-point mutation injected into a short reference solution from
Exercism. Difficulty skews easy, with 66 of 90 tasks (73\%) labeled
\texttt{easy} (single-token or single-operator change) and 24 (27\%) labeled
\texttt{medium}. The functions are short (typically under 20 lines), bugs are
isolated (no cross-function or cross-file effects), and the dataset does not
exercise repository-level reasoning, multi-hunk patches, dependency conflicts,
or subtle bugs requiring long-context understanding. We treat the benchmark as a
controlled testbed for the stem-agent mechanism rather than as a claim about
real-world bug-repair performance, and we report all results with this
constraint in mind. Section~\ref{sec:limitations} discusses this further.

\subsection{Splits and Evaluation Protocol}

% 4 splits: calibration, dev, vital signs, test
% Sizes, stratification, freeze policy
% Metrics: pass@1, bootstrap CI, McNemar
% Baselines: vanilla_direct, vanilla_cot, generic_react, hand_tuned, random_search

\textbf{[TODO Day 2]}: Complete after Phase 0 with actual split sizes.

% =============================================================
% DAY 2: Write Sections 3 + 4.1
% =============================================================

\section{The Stem Agent}

\subsection{Phase 0: Baseline Establishment}

% Baseline results table
% Split assignment procedure

\textbf{[TODO Day 2]}: Fill with Phase 0 results.

\subsection{Phase 1: Sensing}

% What the Sensor analyzes
% Key failure patterns discovered
% Suggested mutations generated

\textbf{[TODO Day 2]}: Fill with Sensor report summary.

\subsection{Phase 2: Differentiation}

% Mutation operators (enumerate key ones)
% Selection mechanism (champion/challenger)
% Safeguard mechanism (vital signs invariant)
% Convergence criterion

\textbf{[TODO Day 2]}: Write mechanism description. \\
\textbf{[TODO Day 3]}: Add differentiation trajectory results.

% =============================================================
% DAY 3: Write Section 4 (Results)
% =============================================================

\section{Results}

\subsection{Main Results}

% Main comparison table: all agents x pass@1 x CI x cost x p-value
% Interpretation

\textbf{[TODO Day 3]}: Fill with test set results.

\begin{table}[h]
\centering
\caption{Test set results (pass@1)}
\begin{tabular}{lcccc}
\toprule
Agent & pass@1 & 95\% CI & Cost/task & $p$ (vs vanilla) \\
\midrule
\texttt{vanilla\_direct} & -- & -- & -- & -- \\
\texttt{vanilla\_cot} & -- & -- & -- & -- \\
\texttt{generic\_react} & -- & -- & -- & -- \\
\texttt{hand\_tuned} & -- & -- & -- & -- \\
\texttt{random\_search} & -- & -- & -- & -- \\
\texttt{stem\_agent} & -- & -- & -- & -- \\
\bottomrule
\end{tabular}
\end{table}

\subsection{Ablation Study}

% Ablation table
% Interpretation of each ablation

\textbf{[TODO Day 3]}: Fill with ablation results.

\subsection{Cross-Domain Validation}

% HumanEval results: stem_agent vs vanilla
% Interpretation

\textbf{[TODO Day 3]}: Fill with cross-domain results.

% =============================================================
% DAY 4: Write Sections 5 + 6 + 7
% =============================================================

\section{Analysis and Discussion}

\subsection{What Worked}

% Most effective mutations
% Categories that improved most
% Role of the Sensor

\textbf{[TODO Day 4]}

\subsection{What Didn't Work}

% Ineffective mutations
% Resistant bug categories
% Search space limitations

\textbf{[TODO Day 4]}

\subsection{The Safeguard in Practice}

% Trigger count
% What it prevented
% Ablation evidence (no_safeguard variant)

\textbf{[TODO Day 4]}

\subsection{Surprises}

% Unexpected findings

\textbf{[TODO Day 4]}

\section{Limitations and Future Work}
\label{sec:limitations}

% Be honest:
% - Small benchmark (N=X)
% - Model mismatch (mini for search, 4o for eval)
% - Hand-designed mutation operators
% - Single domain evaluation
% - Cost of differentiation
% Future:
% - Learned mutations
% - Multi-objective (perf + cost)
% - Transfer between stem agents
% - Larger benchmarks (full SWE-bench)

\textbf{[TODO Day 4]}

\begin{thebibliography}{9}

\bibitem{khattab2023dspy}
Omar Khattab, Arnav Singhvi, Paridhi Maheshwari, Zhiyuan Zhang, Keshav Santhanam, Sri Vardhamanan, Saiful Haq, Ashutosh Sharma, Thomas T. Joshi, Hanna Moazam, Heather Miller, Matei Zaharia, and Christopher Potts.
\newblock {DSPy}: Compiling Declarative Language Model Calls into Self-Improving Pipelines.
\newblock \emph{arXiv preprint arXiv:2310.03714}, 2023.

\bibitem{yang2023opro}
Chengrun Yang, Xuezhi Wang, Yifeng Lu, Hanxiao Liu, Quoc V.\ Le, Denny Zhou, and Xinyun Chen.
\newblock Large Language Models as Optimizers.
\newblock \emph{arXiv preprint arXiv:2309.03409}, 2023.

\bibitem{fernando2023promptbreeder}
Chrisantha Fernando, Dylan Banarse, Henryk Michalewski, Simon Osindero, and Tim Rockt\"{a}schel.
\newblock Promptbreeder: Self-Referential Self-Improvement via Prompt Evolution.
\newblock \emph{arXiv preprint arXiv:2309.16797}, 2023.

\bibitem{hu2024adas}
Shengran Hu, Cong Lu, and Jeff Clune.
\newblock Automated Design of Agentic Systems.
\newblock \emph{arXiv preprint arXiv:2408.08435}, 2024.

\bibitem{jimenez2024swebench}
Carlos E.\ Jim\'{e}nez, John Yang, Alexander Wettig, Shunyu Yao, Kexin Pei, Ofir Press, and Karthik Narasimhan.
\newblock {SWE-bench}: Can Language Models Resolve Real-World {GitHub} Issues?
\newblock In \emph{Proceedings of ICLR}, 2024.

\bibitem{widyasari2020}
Ratnadira Widyasari, Sheng Qin Sim, Camellia Lok, Haodi Qi, Jack Phan, Qijin Tay, Constance Tan, Fiona Wee, Joanna Ethelda Tan, Yuheng Yieh, Brian Goh, Ferdian Thung, Hong Jin Kang, Thiago Syn Wee Soares, Foutse Khomh, and David Lo.
\newblock {BugsInPy}: A database of existing bugs in {Python} programs to enable controlled testing and debugging studies.
\newblock In \emph{Proceedings of ESEC/FSE}, pages 1556--1560, 2020.

\end{thebibliography}

\end{document}